\begin{abstract}
We prove that every bounded Hochschild cocycle of degree at least two
on a complex von Neumann algebra, with values in the algebra itself,
has a bounded primitive. Together with the established degree-one
inner-derivation theorem, this resolves the Kadison--Ringrose cohomology
conjecture positively.
\end{abstract}
\section{Introduction}\label{sec:introduction}

The continuous Hochschild cohomology of an operator algebra measures
the obstruction to solving its bounded cocycle equations. The theory
also connects these equations with stability of multiplication under
small perturbations~\cite[Chapter~7]{SinclairSmith1995}.
For a von Neumann algebra acting on itself by multiplication, the
degree-one equation is the derivation equation, and its vanishing is
the inner-derivation theorem of Kadison and
Sakai~\cite{Kadison1966,Sakai1966}.
The higher-degree theory was developed by Johnson, Kadison, and
Ringrose~\cite{KadisonRingrose1971,KadisonRingrose1971II,JKR1972}.
Throughout, boundedness means ordinary
operator-norm boundedness of a multilinear map. No complete-boundedness
condition is imposed.

For a complex von Neumann algebra \(M\), let \(C_b^k(M,M)\) be the
Banach space of bounded complex \(k\)-linear maps \(M^k\to M\),
with norm
$\norm{f}=\sup_{\norm{a_1},\ldots,\norm{a_k}\leq1}
\norm{f(a_1,\ldots,a_k)}$, and put \(C_b^0(M,M)=M\).
For \(k\geq1\), the differential is
\begin{align}
(df)(a_1,\ldots,a_{k+1})
={}&a_1f(a_2,\ldots,a_{k+1})\notag\\
&+\sum_{j=1}^k(-1)^j
 f(a_1,\ldots,a_ja_{j+1},\ldots,a_{k+1})\notag\\
&+(-1)^{k+1}f(a_1,\ldots,a_k)a_{k+1}.
\label{eq:differential}
\end{align}
In degree zero, \((db)(a)=ab-ba\). As usual, \(d^2=0\).
We use the notation
\[
H_b^k(M,M)=
\ker\bigl(d:C_b^k(M,M)\to C_b^{k+1}(M,M)\bigr)
\big/
\operatorname{im}\bigl(d:C_b^{k-1}(M,M)\to C_b^k(M,M)\bigr).
\]
The denominator is the actual image, without a closure.

The Kadison--Ringrose cohomology conjecture asks whether
\(H_b^k(M,M)=0\) for every complex von Neumann algebra \(M\) and
every positive integer \(k\). Sinclair and Smith trace this
self-coefficient formulation to 1967~\cite[p.~3376]{SinclairSmith1998}.
Our main result settles the higher degrees.

\begin{theorem}\label{thm:main}
Let \(M\) be a complex von Neumann algebra and let \(k\geq2\).
Every \(f\in C_b^k(M,M)\) satisfying \(df=0\) has the form
\(f=dg\) for some \(g\in C_b^{k-1}(M,M)\).
Consequently \(H_b^k(M,M)=0\).
\end{theorem}

Together with the inner-derivation theorem,
Theorem~\ref{thm:main} resolves the Kadison--Ringrose cohomology
conjecture positively.
